\section{LU 7 -- Liệt khuyết (Lieque)}

\subsection{Vị trí huyệt}

Huyệt Liệt khuyết nằm ở mặt quay của cẳng tay, cách nếp gấp ngang cổ tay khoảng \(1{,}5\) thốn về phía gần khuỷu tay, trong khe giữa gân cơ cánh tay quay và gân cơ dạng dài ngón cái.

Có thể xác định huyệt bằng cách đan chéo ngón trỏ và ngón cái của hai bàn tay, đặt ngón trỏ của một bàn tay lên mỏm trâm quay của bàn tay còn lại. Huyệt nằm tại chỗ lõm ngay dưới đầu ngón trỏ.

\begin{figure}[htbp]
    \centering
    \includegraphics[width=0.65\textwidth]{LU_7.jpg}
    \caption{Huyệt Liệt khuyết (LU 7).}
    \label{fig:lu7}
\end{figure}

\subsection{Chỉ định}

Các chỉ định được liệt kê trong tài liệu nguồn bao gồm:

\begin{itemize}
    \item Ho, hen suyễn và đau họng.
    \item Đau đầu, đau nửa đầu và cứng cổ.
    \item Liệt mặt và đau răng.
    \item Đau và yếu vùng cổ tay.
\end{itemize}

\subsection{Cơ chế tác dụng theo y học cổ truyền}

Theo lý luận y học cổ truyền, huyệt Liệt khuyết có các tác dụng sau:

\begin{itemize}
    \item Giải biểu, khu phong.
    \item Tuyên thông Phế khí, thúc đẩy chức năng túc giáng của Phế.
    \item Khai thông và điều hòa mạch Nhâm.
    \item Điều hòa thủy đạo.
    \item Có lợi cho vùng đầu và cổ.
\end{itemize}

\subsection{Ứng dụng trong phương pháp Balance}

Trong phương pháp Balance, huyệt Liệt khuyết được liên hệ với các kinh mạch và hệ thống sau:

\begin{itemize}
    \item \textbf{Kinh thủ Thái âm -- Phế:} kinh mạch chủ quản của huyệt.
    \item \textbf{Hệ thống 1:} kinh túc Thái âm -- Tỳ.
    \item \textbf{Hệ thống 2:} kinh túc Thái dương -- Bàng quang.
    \item \textbf{Hệ thống 3:} kinh thủ Dương minh -- Đại trường.
    \item \textbf{Hệ thống 4:} kinh túc Thái dương -- Bàng quang.
    \item \textbf{Hệ thống 5:} kinh túc Quyết âm -- Can.
\end{itemize}

\subsection{Kỹ thuật châm cứu}

Theo tài liệu nguồn, kỹ thuật châm được mô tả như sau:

\begin{itemize}
    \item \textbf{Hướng châm:} châm chếch hướng lên phía khuỷu tay.
    \item \textbf{Độ sâu:} khoảng \(0{,}3\)--\(0{,}5\) thốn.
\end{itemize}

Hướng châm và độ sâu thực tế cần được điều chỉnh theo đặc điểm giải phẫu, thể trạng người bệnh và đánh giá của người hành nghề có chuyên môn.

\subsection{Đặc tính huyệt}

Theo phân loại trong tài liệu nguồn, Liệt khuyết có các đặc tính sau:

\begin{itemize}
    \item \textbf{Lạc huyệt:} kinh Phế.
    \item \textbf{Huyệt Xuất:} kinh Phế.
    \item \textbf{Huyệt phối:} mạch Âm Kiều (Yinqiao).
    \item \textbf{Huyệt chủ:} mạch Nhâm.
    \item \textbf{Huyệt chỉ huy:} vùng đầu và gáy.
\end{itemize}

\subsection{Tài liệu tham khảo}

Master Tung's Acupuncture. (n.d.). \textit{Lieque (LU 7)}.
eLotus CORE.
\url{https://www.mastertungacupuncture.org/acupuncture/traditional/points/lu7}